\section*{Замечания рецензентов и ответы авторов}
\label{sec:responses}
\noindent Заявка 90. Редакция для DAMDID 2026 от 14 сентября 2026 года.

\paragraph{Рецензент 1: терминология и примеры (скорее принять).}
\emph{Замечание, краткий пересказ русского оригинала.} Статья содержит необходимые разделы исследования, но смысл углублённой доразметки недостаточно ясен. Рецензент просит привести конкретные примеры и пояснения, делающие основные понятия доступными читателям, менее знакомым с темой.

\emph{Ответ.} Во введении и разделе~\ref{sec:annotation} дано операциональное определение: углублённая доразметка расширяет человеческие оценки в отобранном пуле, а не увеличивает глубину нейронной сети и не охватывает весь корпус. Раздел~\ref{sec:example} и таблица~\ref{tab:example} содержат реальный запрос MS MARCO, идентификаторы фрагментов, сокращённые описания их содержания и статус релевантности до и после разметки. Пояснение показывает, почему верный, но неразмеченный ответ не даёт вклада в метрику при неполных qrels. Изменение измеренного качества явно отделено от улучшения неизменных поисковых моделей; объяснены qrels, формирование пула, псевдометки и вход судьи. Эти пояснения включены и в русскую версию.

\paragraph{Рецензент 2: общая оценка (принять).}
\emph{Замечание, краткий пересказ.} Рецензент считает задачу улучшения набора данных важной, результаты ценными, стратегию обоснованной и поддерживает публикацию полноформатной статьи.

\emph{Ответ.} Благодарим рецензента за положительную оценку. Конкретных дополнительных изменений не запрошено. В новой редакции сохранён фокус на охвате разметки релевантности и добавлены изложенные здесь пояснения и методологические уточнения.

\paragraph{Рецензент 3: другие провайдеры и универсальность (уверенно принять).}
\emph{Замечание, краткий пересказ.} Есть ли аналоги MS MARCO у Google, Yandex или Yahoo? Насколько универсальна MS MARCO с учётом возможных различий в использовании поисковых систем между странами и социальными группами? Рецензент предлагает обосновать выбор MS MARCO во введении.

\emph{Ответ.} В разделе~\ref{sec:intro} рассмотрены Natural Questions от Google, Learning to Rank Challenge от Yahoo! и Personalized Web Search Challenge от Yandex со ссылками на первоисточники. Разграничены ответы на вопросы по текстам, ранжирование по векторам признаков и анонимизированные журналы взаимодействий; эти наборы не рассматриваются как взаимозаменяемые. Выбор MS MARCO объяснён доступностью текстов, неполной разметкой, устоявшимися средствами оценки и предобученными моделями. Разделы~\ref{sec:related} и \ref{sec:discussion} обсуждают межколлекционную оценку и прямо указывают, что переносимость между странами и социальными группами не проверялась. Демографическая репрезентативность или универсальная применимость результатов не заявляются.

\paragraph{Дополнительные методологические исправления.}
При проверке установлено, что прежний расчёт использовал отдельный идеальный знаменатель для каждой системы и исключал запросы без найденного положительного фрагмента. Мы перешли к общему IDCG по пулу и единой совокупности запросов, пересчитали таблицы ранжирования, парные эффекты, доверительные интервалы и рисунки. Базовая разметка теперь непосредственно использует официальные qrels подвыборки разработки: в ручном пуле 1\,243 положительные пары; медиана числа ручных положительных оценок исправлена на 17. Псевдометки единообразно восстановлены по исходным оценкам и сохранённому порогу: в ручном пуле 22\,206 положительных пар. Явно указаны вспомогательные источники пула, RRF ограничен тремя названными системами поиска только по запросу, а глубина его входных списков отделена от глубины пула оценки. Эти изменения затрагивают численные результаты и границы выводов, поэтому представлены явно, а не как чисто редакционные правки. Качественный вывод о росте измеренной NDCG после ручной доразметки сохраняет подтверждение при исправленном протоколе.

Направление сравнения объединённой выдачи на полном наборе изменилось относительно прежней версии: восстановленный RRF(3) уступает E5 при глубине десять и превосходит его при четырёх больших исследованных глубинах. Результат сохраняется после исключения всех запросов ручной подвыборки. Раздел~\ref{sec:fullset} и заключение исправлены соответственно; прежнее утверждение о поведении объединения не сохраняется.
